\subsection{Describing Polynomials (Part 1)} 
\begin{definition}
\begin{itemize*}
\item The \term{degree of a polynomial} is the highest degree of any term.
\item The \term{leading term} is the term with highest degree.
\item The \term{leading coefficient} is the coefficient of the leading term.
\end{itemize*}
\noindent If a polynomial has more than one term of the highest degree, it does not have a leading term.
\end{definition}
\begin{Example}
For each polynomial below, identify the leading term, leading coefficient, and degree.
\begin{enumerate}\setabc
\item $3x^3+2x^2+5x+3$
\hfill L. term: \makeblank[0.5in]
\quad L. coeff: \makeblank[0.5in]
\quad Deg: \makeblank[0.5in]
\item $x+\frac{2}{3}x^3-7x^5$
\hfill L. term: \makeblank[0.5in]
\quad L. coeff: \makeblank[0.5in]
\quad Deg: \makeblank[0.5in]
\item $3x^2y-4xy^2+1.3x^2y^2$
\hfill L. term: \makeblank[0.5in]
\quad L. coeff: \makeblank[0.5in]
\quad Deg: \makeblank[0.5in]
\item $x^2-5xy+6y^2$
\hfill L. term: \makeblank[0.5in]
\quad L. coeff: \makeblank[0.5in]
\quad Deg: \makeblank[0.5in]
\end{enumerate}
\end{Example}
\begin{definition}
\begin{itemize*}
\item A \term{polynomial in one variable} is one in which only one variable appears.
\item A \term{homogeneous} polynomial is one in which every term has equal degree.
\end{itemize*}
\end{definition}
\begin{Example}
For each polynomial below, identify it as a polynomial in one variable, a homogeneous polynomial, or neither. Also determine the degree of the polynomial.
\begin{enumerate}\setabc
\item $x^4-x^2+2x+1$
\hfill One variable {\textbar} Homogeneous {\textbar} Neither
\qquad Degree: \makeblanksmall
\item $x^2-5xy-6y^2$
\hfill One variable {\textbar} Homogeneous {\textbar} Neither
\qquad Degree: \makeblanksmall
\item $6xy-10x+9y$
\hfill One variable {\textbar} Homogeneous {\textbar} Neither
\qquad Degree: \makeblanksmall
\end{enumerate}
\end{Example}
\begin{definition}
\begin{itemize*}
\item A polynomial is written in \term{descending order} if the terms are ordered from highest to lowest degree.
\item A polynomial is written in \term{ascending order} if the terms are ordered from lowest to highest degree.
\end{itemize*}
\end{definition}
\begin{Example}
Write $5x-3x^2-2x^3+6$ in descending order.\makeblanknewf
\end{Example}
\begin{Example}
Write $5y+15-6y^5+2y^7$ in ascending order.\makeblanknewf
\end{Example}